%% ma: L12-B06-0003
%% lop: 12
%% chuong: 2
%% bai: 6
%% dang: 12.6.3
%% loai: TN4
%% muc: TH
%% dap_an: "B"
%% nguon: "(NBV)-ĐỀ SỐ 3-MỖI NGÀY 1 ĐỀ THI 2025.pdf (D03, MỖI NGÀY 1 ĐỀ - Đề số 3), Phần I Câu 4"
%% kien_thuc: "Bài 6 (tổng hai vectơ, trung điểm đoạn thẳng, vectơ trong không gian)"
%% trang_thai: cho-duyet
%% ghi_chu: "Đáp án gốc B đúng. Hình dựng lại: BD là nét khuất vì C là đỉnh gần người nhìn"
%% ly_do: "Dạng toán mới đề xuất 12.6.3 chờ giáo viên duyệt"
\begin{ex}
Cho tứ diện $ABCD$. Gọi $I,J,O$ theo thứ tự là trung điểm của các đoạn thẳng $AB,CD,IJ$ như hình vẽ. Khẳng định nào sau đây sai?
\begin{center}
\begin{tikzpicture}
  \coordinate (A) at (2.4,3.4);
  \coordinate (B) at (0,0.5);
  \coordinate (C) at (1.0,-1.3);
  \coordinate (D) at (4.6,0);
  \coordinate (I) at ($(A)!.5!(B)$);
  \coordinate (J) at ($(C)!.5!(D)$);
  \coordinate (O) at ($(I)!.5!(J)$);
  \draw (A)--(B)--(C)--(D)--cycle;
  \draw (A)--(C);
  \draw[khuat] (B)--(D);
  \draw[khuat] (I)--(J);
  \foreach \p in {A,B,C,D,I,J,O} \fill (\p) circle (1.2pt);
  \node[above] at (A) {$A$};
  \node[left] at (B) {$B$};
  \node[below] at (C) {$C$};
  \node[right] at (D) {$D$};
  \node[above left] at (I) {$I$};
  \node[below right] at (J) {$J$};
  \node[above right] at (O) {$O$};
\end{tikzpicture}
\end{center}
\choice{$\overrightarrow{IA}+\overrightarrow{IB}=\vec0$}{\True $\overrightarrow{AI}+\overrightarrow{AJ}=\dfrac12\overrightarrow{AO}$}{$\overrightarrow{OC}+\overrightarrow{OD}=2\overrightarrow{OJ}$}{$\overrightarrow{OJ}=\dfrac12\left(\overrightarrow{OC}+\overrightarrow{OD}\right)$}
\loigiai{$I$ là trung điểm $AB$ nên $\overrightarrow{IA}+\overrightarrow{IB}=\vec0$; $J$ là trung điểm $CD$ nên $\overrightarrow{OC}+\overrightarrow{OD}=2\overrightarrow{OJ}$, suy ra các khẳng định A, C, D đúng. $O$ là trung điểm $IJ$ nên $\overrightarrow{AI}+\overrightarrow{AJ}=2\overrightarrow{AO}\neq\dfrac12\overrightarrow{AO}$. Vậy khẳng định B sai.}
\end{ex}
